\SGASection{2}{Reductive group schemes: definitions and first properties}
\SGARefTarget{sga3:xix:section:2}

\par\medskip\noindent
\textbf{\textit{\SGARefLink{sga3:xix:scholium:2:1}{Scholium 2.1}.}}---
\textit{For an \(S\)-group scheme \(G\), the following properties are
equivalent:
\begin{enumerate}
\item[\SGAInlineRefTarget{sga3:xix:section:2:item:i}(i)] \(G\) is affine and smooth over \(S\), with connected
  fibers.
\item[\SGAInlineRefTarget{sga3:xix:section:2:item:ii}(ii)] \(G\) is affine and flat over \(S\), of finite presentation
  over \(S\), with geometrically integral fibers.
\end{enumerate}
These properties are stable under base change and are local for the
fpqc topology.}
\SGARefTarget{sga3:xix:scholium:2:1}

\footnote{Editor's note: the following details have been added.}
Indeed, suppose that \SGARefLink{sga3:xix:section:2:item:i}{(i)} holds.  Since \(G\) is affine and smooth over
\(S\), it is of finite presentation over \(S\); and since its fibers
are smooth and connected, they are geometrically integral, by
Exposé VI\textsubscript{A}, \SGARefLink{prop:sga3-VIA-2.4}{2.4}.

Conversely, if \SGARefLink{sga3:xix:section:2:item:ii}{(ii)} holds, the fibers of \(G\) are connected and
geometrically reduced, and hence smooth
(Exposé VI\textsubscript{A}, 1.3.1).  It follows that \(G\) is smooth
over \(S\), by EGA IV\textsubscript{4}, 17.5.2.

These properties are plainly stable under base change: see EGA II,
1.5.1 for ``affine'', EGA IV\textsubscript{1}, 1.6.2(iii) for ``of
finite presentation'', EGA IV\textsubscript{2}, 2.1.4 for ``flat'',
and EGA IV\textsubscript{2}, 4.6.5(i) for ``with geometrically integral
fibers''.

They are also clearly local for the Zariski topology.  It is therefore
enough to check that, if \(S'\longrightarrow S\) is faithfully flat and
quasi-compact and \(G_{S'}\longrightarrow S'\) has the indicated
properties, then so does \(G\longrightarrow S\).  This follows from
EGA IV\textsubscript{2}, 2.5.1 for ``flat'', 2.7.1(vi) and (xiii) for
``of finite presentation'' and ``affine'', and 4.6.5(i) for ``with
geometrically integral fibers''; for ``smooth'', see also
EGA IV\textsubscript{4}, 17.7.3(ii).

\par\medskip\noindent\SGAInlineRefTarget{sga3:xix:paragraph:2-2}\textbf{2.2.}\quad
Let \(G\) be an \(S\)-group scheme satisfying the preceding conditions,
and let \(Q\) be a torus of \(G\) in the sense of \SGARefLink{sga3:ix:definition:1:3}{Exposé IX,
Definition 1.3}.\footnote{Editor's note: the following argument has been
given in detail.}
By \SGARefLink{sga3:xi:corollary:6:11:item:a}{Exposé XI, 6.11(a)} and \SGARefLink{sga3:xix:remark:2-4}{2.4},
\(\underline{\operatorname{Centr}}_{G}(Q)\) is represented by a closed
subgroup scheme of \(G\); it is therefore affine over \(S\), and it is
of finite presentation and smooth over \(S\).  Moreover, every
geometric fiber \(G_{\bar s}\) is a smooth affine connected
\(\kappa(\bar s)\)-group.  The first assertion of \SGARefLink{sga3:xix:paragraph:1-3}{1.3} therefore shows
that the same is true of
\[
  \underline{\operatorname{Centr}}_{G_{\bar s}}(Q_{\bar s})
  =
  \bigl(\underline{\operatorname{Centr}}_{G}(Q)\bigr)_{\bar s}.
\]
\SGARefTarget{sga3:xix:section:2:2}
